\chapter{Unit 5: Workflow Design and Iterative Refinement}


\section{Section A: Multiple Choice Practice Questions (60 Curated Problems)}

\noindent Practice and test your mastery across all core concepts of this unit. Each question reflects standard university examination patterns and Bloom's Taxonomy competencies.
\vspace{3mm}

\noindent\textbf{Q1. What is a workflow?} \hfill [Tag: Introduction to workflow and process design for solution development | Easy]
\begin{enumerate}[label=(\alph*),topsep=2pt,itemsep=1pt,partopsep=0pt,parsep=0pt]
    \item {A method for avoiding all decisions}
    \item {A final product without instructions}
    \item {A sequence of steps for completing a task}
    \item {A list of unrelated ideas}
\end{enumerate}
\vspace{2mm}

\noindent\textbf{Q2. What is the main purpose of process design?} \hfill [Tag: Introduction to workflow and process design for solution development | Easy]
\begin{enumerate}[label=(\alph*),topsep=2pt,itemsep=1pt,partopsep=0pt,parsep=0pt]
    \item {To remove all user involvement}
    \item {To make every task identical}
    \item {To delay decisions until the end}
    \item {To organize activities toward a goal}
\end{enumerate}
\vspace{2mm}

\noindent\textbf{Q3. Which tool commonly represents the steps in a workflow?} \hfill [Tag: Introduction to workflow and process design for solution development | Easy]
\begin{enumerate}[label=(\alph*),topsep=2pt,itemsep=1pt,partopsep=0pt,parsep=0pt]
    \item {Glossary}
    \item {Budget sheet}
    \item {Flowchart}
    \item {Survey scale}
\end{enumerate}
\vspace{2mm}

\noindent\textbf{Q4. What does a decision point in a workflow show?} \hfill [Tag: Introduction to workflow and process design for solution development | Easy]
\begin{enumerate}[label=(\alph*),topsep=2pt,itemsep=1pt,partopsep=0pt,parsep=0pt]
    \item {A list of team members}
    \item {A choice between possible paths}
    \item {A summary of project costs}
    \item {A record of completed meetings}
\end{enumerate}
\vspace{2mm}

\noindent\textbf{Q5. What is validation?} \hfill [Tag: Methods for validation | Easy]
\begin{enumerate}[label=(\alph*),topsep=2pt,itemsep=1pt,partopsep=0pt,parsep=0pt]
    \item {Writing instructions before defining a problem}
    \item {Adding features without research}
    \item {Choosing a color for a document}
    \item {Checking whether a solution meets user needs}
\end{enumerate}
\vspace{2mm}

\noindent\textbf{Q6. Which activity is a common method of validating an idea?} \hfill [Tag: Methods for validation | Easy]
\begin{enumerate}[label=(\alph*),topsep=2pt,itemsep=1pt,partopsep=0pt,parsep=0pt]
    \item {Skipping all observations}
    \item {Gathering user feedback}
    \item {Ignoring user comments}
    \item {Deleting the project plan}
\end{enumerate}
\vspace{2mm}

\noindent\textbf{Q7. What is a prototype used for during validation?} \hfill [Tag: Methods for validation | Easy]
\begin{enumerate}[label=(\alph*),topsep=2pt,itemsep=1pt,partopsep=0pt,parsep=0pt]
    \item {Preventing users from seeing the idea}
    \item {Measuring only the project budget}
    \item {Exploring and evaluating an idea}
    \item {Replacing every future product}
\end{enumerate}
\vspace{2mm}

\noindent\textbf{Q8. Which method asks users questions to learn about their needs?} \hfill [Tag: Methods for validation | Easy]
\begin{enumerate}[label=(\alph*),topsep=2pt,itemsep=1pt,partopsep=0pt,parsep=0pt]
    \item {User interview}
    \item {File compression}
    \item {Task deletion}
    \item {Color coding}
\end{enumerate}
\vspace{2mm}

\noindent\textbf{Q9. What is the main purpose of testing a solution?} \hfill [Tag: Testing | Easy]
\begin{enumerate}[label=(\alph*),topsep=2pt,itemsep=1pt,partopsep=0pt,parsep=0pt]
    \item {To guarantee that no changes are needed}
    \item {To replace the need for users}
    \item {To make the workflow longer}
    \item {To find problems and check performance}
\end{enumerate}
\vspace{2mm}

\noindent\textbf{Q10. What is usability testing focused on?} \hfill [Tag: Testing | Easy]
\begin{enumerate}[label=(\alph*),topsep=2pt,itemsep=1pt,partopsep=0pt,parsep=0pt]
    \item {How quickly a team holds meetings}
    \item {How easily users can use a solution}
    \item {How much office space is available}
    \item {How many pages a report contains}
\end{enumerate}
\vspace{2mm}

\noindent\textbf{Q11. When should testing usually begin in an iterative process?} \hfill [Tag: Testing | Easy]
\begin{enumerate}[label=(\alph*),topsep=2pt,itemsep=1pt,partopsep=0pt,parsep=0pt]
    \item {When the problem is forgotten}
    \item {Only after the final launch}
    \item {After all feedback is discarded}
    \item {As early as practical}
\end{enumerate}
\vspace{2mm}

\noindent\textbf{Q12. What does a test result provide?} \hfill [Tag: Testing | Easy]
\begin{enumerate}[label=(\alph*),topsep=2pt,itemsep=1pt,partopsep=0pt,parsep=0pt]
    \item {Evidence about how a solution works}
    \item {A list of unrelated solutions}
    \item {A replacement for the original problem}
    \item {A guarantee of permanent success}
\end{enumerate}
\vspace{2mm}

\noindent\textbf{Q13. What does iterative refinement mean?} \hfill [Tag: Iterative refinement | Easy]
\begin{enumerate}[label=(\alph*),topsep=2pt,itemsep=1pt,partopsep=0pt,parsep=0pt]
    \item {Completing a solution without review}
    \item {Improving a solution through repeated changes}
    \item {Choosing the first idea immediately}
    \item {Removing all versions of a solution}
\end{enumerate}
\vspace{2mm}

\noindent\textbf{Q14. What usually happens after a solution is tested in an iterative process?} \hfill [Tag: Iterative refinement | Easy]
\begin{enumerate}[label=(\alph*),topsep=2pt,itemsep=1pt,partopsep=0pt,parsep=0pt]
    \item {All test results are discarded}
    \item {The users are removed from the process}
    \item {The problem is automatically solved forever}
    \item {The solution is reviewed and improved}
\end{enumerate}
\vspace{2mm}

\noindent\textbf{Q15. Why are multiple versions of a solution useful?} \hfill [Tag: Iterative refinement | Easy]
\begin{enumerate}[label=(\alph*),topsep=2pt,itemsep=1pt,partopsep=0pt,parsep=0pt]
    \item {They prevent comparison between ideas}
    \item {They eliminate the need for testing}
    \item {They allow learning and gradual improvement}
    \item {They ensure that the first version remains unchanged}
\end{enumerate}
\vspace{2mm}

\noindent\textbf{Q16. What is feedback?} \hfill [Tag: Feedback incorporation | Easy]
\begin{enumerate}[label=(\alph*),topsep=2pt,itemsep=1pt,partopsep=0pt,parsep=0pt]
    \item {Information about a user's experience or response}
    \item {A fixed schedule for team holidays}
    \item {A final answer that cannot change}
    \item {A list of materials used in construction}
\end{enumerate}
\vspace{2mm}

\noindent\textbf{Q17. How should useful feedback be handled?} \hfill [Tag: Feedback incorporation | Easy]
\begin{enumerate}[label=(\alph*),topsep=2pt,itemsep=1pt,partopsep=0pt,parsep=0pt]
    \item {Ignore it until the project ends}
    \item {Review it and use it to guide changes}
    \item {Accept every suggestion without review}
    \item {Replace testing with personal opinions}
\end{enumerate}
\vspace{2mm}

\noindent\textbf{Q18. Which source can provide feedback on a proposed solution?} \hfill [Tag: Feedback incorporation | Easy]
\begin{enumerate}[label=(\alph*),topsep=2pt,itemsep=1pt,partopsep=0pt,parsep=0pt]
    \item {Random page numbers}
    \item {Intended users}
    \item {Unrelated advertisements}
    \item {Unused file names}
\end{enumerate}
\vspace{2mm}

\noindent\textbf{Q19. What is continuous improvement?} \hfill [Tag: Continuous improvement of proposed solutions | Easy]
\begin{enumerate}[label=(\alph*),topsep=2pt,itemsep=1pt,partopsep=0pt,parsep=0pt]
    \item {Stopping development after one attempt}
    \item {Avoiding all evaluation activities}
    \item {Changing goals without collecting information}
    \item {Making ongoing changes to improve results}
\end{enumerate}
\vspace{2mm}

\noindent\textbf{Q20. Which activity supports continuous improvement?} \hfill [Tag: Continuous improvement of proposed solutions | Easy]
\begin{enumerate}[label=(\alph*),topsep=2pt,itemsep=1pt,partopsep=0pt,parsep=0pt]
    \item {Ending feedback collection immediately}
    \item {Keeping known problems unchanged}
    \item {Regularly reviewing performance}
    \item {Ignoring new evidence}
\end{enumerate}
\vspace{2mm}

\noindent\textbf{Q21. A team is designing a mobile service to help patients schedule appointments. Which workflow representation would best clarify the sequence of user actions and system responses?} \hfill [Tag: Introduction to workflow and process design for solution development | Medium]
\begin{enumerate}[label=(\alph*),topsep=2pt,itemsep=1pt,partopsep=0pt,parsep=0pt]
    \item {A user journey flowchart}
    \item {A list of technical features}
    \item {A project cost summary}
    \item {A stakeholder influence map}
\end{enumerate}
\vspace{2mm}

\noindent\textbf{Q22. A workflow contains five approval stages, but most requests require only two. What is the most appropriate initial design response?} \hfill [Tag: Introduction to workflow and process design for solution development | Medium]
\begin{enumerate}[label=(\alph*),topsep=2pt,itemsep=1pt,partopsep=0pt,parsep=0pt]
    \item {Add another approval stage}
    \item {Allow every request to bypass review}
    \item {Identify which stages add necessary value}
    \item {Remove all approval stages}
\end{enumerate}
\vspace{2mm}

\noindent\textbf{Q23. In a solution-development workflow, what is the main purpose of defining inputs and outputs for each activity?} \hfill [Tag: Introduction to workflow and process design for solution development | Medium]
\begin{enumerate}[label=(\alph*),topsep=2pt,itemsep=1pt,partopsep=0pt,parsep=0pt]
    \item {To assign blame for delays}
    \item {To increase the number of workflow steps}
    \item {To clarify dependencies between activities}
    \item {To replace user research with documentation}
\end{enumerate}
\vspace{2mm}

\noindent\textbf{Q24. A team notices that work frequently waits between research and prototyping. Which workflow improvement is most likely to address the problem?} \hfill [Tag: Introduction to workflow and process design for solution development | Medium]
\begin{enumerate}[label=(\alph*),topsep=2pt,itemsep=1pt,partopsep=0pt,parsep=0pt]
    \item {Measure only the team's total hours}
    \item {Create a clearer handoff agreement}
    \item {Delay prototyping until final approval}
    \item {Increase the number of research questions}
\end{enumerate}
\vspace{2mm}

\noindent\textbf{Q25. A team wants to determine whether users can understand a new dashboard without assistance. Which validation method is most suitable?} \hfill [Tag: Methods for validation | Medium]
\begin{enumerate}[label=(\alph*),topsep=2pt,itemsep=1pt,partopsep=0pt,parsep=0pt]
    \item {A competitor feature comparison}
    \item {A moderated usability test}
    \item {A technical code inspection}
    \item {A financial feasibility review}
\end{enumerate}
\vspace{2mm}

\noindent\textbf{Q26. A prototype receives positive comments from five team members, but none are target users. What is the main limitation of this validation?} \hfill [Tag: Methods for validation | Medium]
\begin{enumerate}[label=(\alph*),topsep=2pt,itemsep=1pt,partopsep=0pt,parsep=0pt]
    \item {The team members used different devices}
    \item {The feedback lacks target-user relevance}
    \item {The comments were collected too quickly}
    \item {The prototype has too many screens}
\end{enumerate}
\vspace{2mm}

\noindent\textbf{Q27. A team is choosing between two onboarding designs and wants evidence about which one leads to more completed registrations. Which method is most appropriate?} \hfill [Tag: Methods for validation | Medium]
\begin{enumerate}[label=(\alph*),topsep=2pt,itemsep=1pt,partopsep=0pt,parsep=0pt]
    \item {A brainstorm of additional features}
    \item {A/B testing with a defined success metric}
    \item {An informal discussion with designers}
    \item {A review of the project schedule}
\end{enumerate}
\vspace{2mm}

\noindent\textbf{Q28. A proposed solution appears desirable but requires expensive infrastructure. Which validation method should be prioritized before full development?} \hfill [Tag: Methods for validation | Medium]
\begin{enumerate}[label=(\alph*),topsep=2pt,itemsep=1pt,partopsep=0pt,parsep=0pt]
    \item {A color preference survey}
    \item {A final launch campaign}
    \item {A feasibility assessment}
    \item {A retrospective meeting}
\end{enumerate}
\vspace{2mm}

\noindent\textbf{Q29. A survey shows that 80\% of respondents like a proposed feature, but interviews reveal that many would not use it regularly. What should the team conclude?} \hfill [Tag: Methods for validation | Medium]
\begin{enumerate}[label=(\alph*),topsep=2pt,itemsep=1pt,partopsep=0pt,parsep=0pt]
    \item {The interviews should be discarded as subjective}
    \item {The feature should be launched unchanged}
    \item {The survey alone proves adoption}
    \item {The evidence should be examined using multiple methods}
\end{enumerate}
\vspace{2mm}

\noindent\textbf{Q30. A team tests a prototype with a fictional user profile that does not represent the intended audience. What is the most likely consequence?} \hfill [Tag: Testing | Medium]
\begin{enumerate}[label=(\alph*),topsep=2pt,itemsep=1pt,partopsep=0pt,parsep=0pt]
    \item {The workflow will require fewer decisions}
    \item {The findings may not generalize to real users}
    \item {The test will measure development speed more accurately}
    \item {The prototype will automatically become more reliable}
\end{enumerate}
\vspace{2mm}

\noindent\textbf{Q31. During a usability test, nearly every participant fails at the same navigation step. What should the team do first?} \hfill [Tag: Testing | Medium]
\begin{enumerate}[label=(\alph*),topsep=2pt,itemsep=1pt,partopsep=0pt,parsep=0pt]
    \item {Redesign the navigation hypothesis}
    \item {Recruit only expert users}
    \item {Increase the product price}
    \item {Ignore the result until launch}
\end{enumerate}
\vspace{2mm}

\noindent\textbf{Q32. Why should a testing team define success criteria before running a test?} \hfill [Tag: Testing | Medium]
\begin{enumerate}[label=(\alph*),topsep=2pt,itemsep=1pt,partopsep=0pt,parsep=0pt]
    \item {To prevent all unexpected findings}
    \item {To guarantee positive participant responses}
    \item {To eliminate the need for observation}
    \item {To make results easier to interpret}
\end{enumerate}
\vspace{2mm}

\noindent\textbf{Q33. A prototype passes functional checks but users take twice as long as expected to complete the main task. What does this result indicate?} \hfill [Tag: Testing | Medium]
\begin{enumerate}[label=(\alph*),topsep=2pt,itemsep=1pt,partopsep=0pt,parsep=0pt]
    \item {The functional checks were unnecessary}
    \item {The solution has no remaining design risks}
    \item {The solution is technically functional but needs usability improvement}
    \item {The users should be removed from future tests}
\end{enumerate}
\vspace{2mm}

\noindent\textbf{Q34. A team has limited time and identifies ten possible improvements. Which approach best supports iterative refinement?} \hfill [Tag: Iterative refinement | Medium]
\begin{enumerate}[label=(\alph*),topsep=2pt,itemsep=1pt,partopsep=0pt,parsep=0pt]
    \item {Prioritize changes by impact and effort}
    \item {Choose only changes requested by managers}
    \item {Implement changes in random order}
    \item {Make every change before testing again}
\end{enumerate}
\vspace{2mm}

\noindent\textbf{Q35. A team changes the interface, pricing, and target audience at the same time before retesting. Why is this approach problematic?} \hfill [Tag: Iterative refinement | Medium]
\begin{enumerate}[label=(\alph*),topsep=2pt,itemsep=1pt,partopsep=0pt,parsep=0pt]
    \item {It becomes difficult to identify which change affected results}
    \item {It makes the prototype too visually simple}
    \item {It guarantees that the next test will fail}
    \item {It prevents stakeholders from seeing the prototype}
\end{enumerate}
\vspace{2mm}

\noindent\textbf{Q36. After each prototype cycle, a team records the change made, the evidence supporting it, and the observed result. What is the main benefit?} \hfill [Tag: Iterative refinement | Medium]
\begin{enumerate}[label=(\alph*),topsep=2pt,itemsep=1pt,partopsep=0pt,parsep=0pt]
    \item {It reduces the solution to one feature}
    \item {It removes the need for future testing}
    \item {It ensures every stakeholder agrees}
    \item {It creates traceability for design decisions}
\end{enumerate}
\vspace{2mm}

\noindent\textbf{Q37. Users suggest adding many features, but several suggestions conflict with the primary goal of the product. How should the team respond?} \hfill [Tag: Feedback incorporation | Medium]
\begin{enumerate}[label=(\alph*),topsep=2pt,itemsep=1pt,partopsep=0pt,parsep=0pt]
    \item {Evaluate suggestions against the product goal}
    \item {Add every suggestion immediately}
    \item {Reject all suggestions without review}
    \item {Select the suggestions from the largest group}
\end{enumerate}
\vspace{2mm}

\noindent\textbf{Q38. A stakeholder requests a change based on one unusual incident, while repeated user tests show a different pattern. What is the best response?} \hfill [Tag: Feedback incorporation | Medium]
\begin{enumerate}[label=(\alph*),topsep=2pt,itemsep=1pt,partopsep=0pt,parsep=0pt]
    \item {Ignore the stakeholder request completely}
    \item {Compare the incident with the broader evidence}
    \item {Cancel all further testing activities}
    \item {Adopt the stakeholder request automatically}
\end{enumerate}
\vspace{2mm}

\noindent\textbf{Q39. Which practice helps prevent feedback from becoming a list of disconnected opinions?} \hfill [Tag: Feedback incorporation | Medium]
\begin{enumerate}[label=(\alph*),topsep=2pt,itemsep=1pt,partopsep=0pt,parsep=0pt]
    \item {Remove feedback that suggests major changes}
    \item {Group feedback by themes and user needs}
    \item {Prioritize comments by emotional intensity}
    \item {Record only comments from senior staff}
\end{enumerate}
\vspace{2mm}

\noindent\textbf{Q40. A service meets its launch targets, but monthly data shows a gradual decline in successful task completion. What should the team do?} \hfill [Tag: Continuous improvement of proposed solutions | Medium]
\begin{enumerate}[label=(\alph*),topsep=2pt,itemsep=1pt,partopsep=0pt,parsep=0pt]
    \item {Treat the launch as the final evaluation}
    \item {Investigate the decline and plan another improvement cycle}
    \item {Replace the service without examining evidence}
    \item {Remove the success metric from future reports}
\end{enumerate}
\vspace{2mm}

\noindent\textbf{Q41. A team is designing a workflow for resolving customer billing disputes. Cases vary greatly in complexity, but every case must receive a compliance review. Which workflow structure best balances efficiency and control?} \hfill [Tag: Introduction to workflow and process design for solution development | Hard]
\begin{enumerate}[label=(\alph*),topsep=2pt,itemsep=1pt,partopsep=0pt,parsep=0pt]
    \item {Allow agents to choose any sequence based on judgment}
    \item {Use one identical sequence for every dispute}
    \item {Route all disputes directly to compliance}
    \item {Use risk-based branching with a mandatory compliance checkpoint}
\end{enumerate}
\vspace{2mm}

\noindent\textbf{Q42. A proposed workflow reduces average completion time but increases handoffs from three to nine. Which analysis is most important before adoption?} \hfill [Tag: Introduction to workflow and process design for solution development | Hard]
\begin{enumerate}[label=(\alph*),topsep=2pt,itemsep=1pt,partopsep=0pt,parsep=0pt]
    \item {Estimate handoff failure costs and queue effects}
    \item {Ask users whether the diagram looks simpler}
    \item {Compare only the new average completion time}
    \item {Select the workflow with fewer documented steps}
\end{enumerate}
\vspace{2mm}

\noindent\textbf{Q43. In a workflow, activities A and B can occur in parallel, but activity C requires both outputs. If A takes 4 minutes and B takes 7 minutes, excluding coordination overhead, what is the minimum duration from starting A and B to completing C?} \hfill [Tag: Introduction to workflow and process design for solution development | Hard]
\begin{enumerate}[label=(\alph*),topsep=2pt,itemsep=1pt,partopsep=0pt,parsep=0pt]
    \item {11 minutes}
    \item {4 minutes}
    \item {7 minutes}
    \item {28 minutes}
\end{enumerate}
\vspace{2mm}

\noindent\textbf{Q44. A service workflow has a high rate of rework because users submit incomplete information. Which redesign is most likely to address the root cause rather than its symptoms?} \hfill [Tag: Introduction to workflow and process design for solution development | Hard]
\begin{enumerate}[label=(\alph*),topsep=2pt,itemsep=1pt,partopsep=0pt,parsep=0pt]
    \item {Add a final audit before service delivery}
    \item {Increase staffing in the rework queue}
    \item {Measure rework duration more frequently}
    \item {Add early validation and contextual input guidance}
\end{enumerate}
\vspace{2mm}

\noindent\textbf{Q45. A prototype claims to reduce medication errors. Which validation design provides the strongest evidence while limiting confounding from differences between clinical units?} \hfill [Tag: Methods for validation | Hard]
\begin{enumerate}[label=(\alph*),topsep=2pt,itemsep=1pt,partopsep=0pt,parsep=0pt]
    \item {Measure how quickly clinicians complete the prototype tasks}
    \item {Ask clinicians to rate the prototype after a demonstration}
    \item {Randomize comparable cases to prototype and standard workflows}
    \item {Compare outcomes before and after deployment in one unit}
\end{enumerate}
\vspace{2mm}

\noindent\textbf{Q46. A validation study reports a statistically significant improvement, but the effect is smaller than the minimum benefit stakeholders consider worthwhile. What is the best interpretation?} \hfill [Tag: Methods for validation | Hard]
\begin{enumerate}[label=(\alph*),topsep=2pt,itemsep=1pt,partopsep=0pt,parsep=0pt]
    \item {The result is invalid because small effects cannot be significant}
    \item {The solution is valuable because significance proves usefulness}
    \item {The solution may be reliable but not practically valuable}
    \item {The sample should be enlarged until the benefit becomes larger}
\end{enumerate}
\vspace{2mm}

\noindent\textbf{Q47. A team validates a workflow only with expert users and observes excellent performance. Which threat remains most serious?} \hfill [Tag: Methods for validation | Hard]
\begin{enumerate}[label=(\alph*),topsep=2pt,itemsep=1pt,partopsep=0pt,parsep=0pt]
    \item {Expert feedback cannot identify any usability issue}
    \item {Experts cannot complete complex workflow tasks}
    \item {Experts always prefer slower workflows}
    \item {Experts may produce an artificially low error rate}
\end{enumerate}
\vspace{2mm}

\noindent\textbf{Q48. A team wants to validate whether a proposed intake process serves users with intermittent internet access. Which method is most appropriate initially?} \hfill [Tag: Methods for validation | Hard]
\begin{enumerate}[label=(\alph*),topsep=2pt,itemsep=1pt,partopsep=0pt,parsep=0pt]
    \item {Replace the process with an offline-only design}
    \item {Use scenario-based trials with controlled interruptions}
    \item {Test only under ideal network conditions}
    \item {Infer performance from desktop usability ratings}
\end{enumerate}
\vspace{2mm}

\noindent\textbf{Q49. A workflow passes every individual component test but fails when two automated services update the same record simultaneously. Which testing gap does this reveal?} \hfill [Tag: Testing | Hard]
\begin{enumerate}[label=(\alph*),topsep=2pt,itemsep=1pt,partopsep=0pt,parsep=0pt]
    \item {Insufficient unit testing of each service}
    \item {Insufficient integration and concurrency testing}
    \item {Insufficient collection of stakeholder preferences}
    \item {Insufficient visual inspection of the workflow}
\end{enumerate}
\vspace{2mm}

\noindent\textbf{Q50. A test suite shows 98\% pass rate, but all failed tests concern a low-frequency safety exception. What should the team do first?} \hfill [Tag: Testing | Hard]
\begin{enumerate}[label=(\alph*),topsep=2pt,itemsep=1pt,partopsep=0pt,parsep=0pt]
    \item {Average the failures across all workflow stages}
    \item {Ignore the failures because the overall rate is high}
    \item {Prioritize the exception based on consequence severity}
    \item {Remove the exception from the acceptance criteria}
\end{enumerate}
\vspace{2mm}

\noindent\textbf{Q51. A workflow test produces inconsistent results because an external data service changes during each run. Which improvement most directly increases test reliability?} \hfill [Tag: Testing | Hard]
\begin{enumerate}[label=(\alph*),topsep=2pt,itemsep=1pt,partopsep=0pt,parsep=0pt]
    \item {Replace outcome measures with user opinions}
    \item {Run the same test more times without controls}
    \item {Use a versioned test double with fixed responses}
    \item {Exclude all cases requiring external data}
\end{enumerate}
\vspace{2mm}

\noindent\textbf{Q52. A new workflow improves average throughput but causes severe delays for the smallest 5\% of cases. Which test metric best exposes this problem?} \hfill [Tag: Testing | Hard]
\begin{enumerate}[label=(\alph*),topsep=2pt,itemsep=1pt,partopsep=0pt,parsep=0pt]
    \item {Overall number of completed cases}
    \item {Tail latency such as the 95th percentile}
    \item {Median completion time only}
    \item {Average number of workflow activities}
\end{enumerate}
\vspace{2mm}

\noindent\textbf{Q53. During iterative refinement, users request five changes, but engineering capacity permits only two. Which selection rule is most defensible?} \hfill [Tag: Iterative refinement | Hard]
\begin{enumerate}[label=(\alph*),topsep=2pt,itemsep=1pt,partopsep=0pt,parsep=0pt]
    \item {Choose the requests requiring the least development effort}
    \item {Choose changes with the highest validated impact-to-cost ratio}
    \item {Choose the requests from the most senior users}
    \item {Choose the requests mentioned most recently}
\end{enumerate}
\vspace{2mm}

\noindent\textbf{Q54. A team makes several simultaneous workflow changes and observes improved outcomes. Why is this iteration difficult to learn from?} \hfill [Tag: Iterative refinement | Hard]
\begin{enumerate}[label=(\alph*),topsep=2pt,itemsep=1pt,partopsep=0pt,parsep=0pt]
    \item {Multiple changes make the workflow too short}
    \item {The team cannot attribute effects to individual changes}
    \item {Simultaneous changes prevent stakeholder participation}
    \item {Improved outcomes always indicate testing bias}
\end{enumerate}
\vspace{2mm}

\noindent\textbf{Q55. A refined solution performs better in a pilot but worsens after full deployment because users develop workarounds. Which refinement practice would best reduce this risk?} \hfill [Tag: Iterative refinement | Hard]
\begin{enumerate}[label=(\alph*),topsep=2pt,itemsep=1pt,partopsep=0pt,parsep=0pt]
    \item {Study real-world use and test likely adaptations}
    \item {Measure only the intended official sequence}
    \item {Freeze the design immediately after pilot success}
    \item {Remove all customization from the workflow}
\end{enumerate}
\vspace{2mm}

\noindent\textbf{Q56. Users report that a workflow is confusing, but their proposed fix would increase compliance risk. What should the design team do?} \hfill [Tag: Feedback incorporation | Hard]
\begin{enumerate}[label=(\alph*),topsep=2pt,itemsep=1pt,partopsep=0pt,parsep=0pt]
    \item {Keep the workflow unchanged to preserve compliance}
    \item {Identify the underlying need and redesign within constraints}
    \item {Implement the proposed fix without modification}
    \item {Reject the feedback because users are not designers}
\end{enumerate}
\vspace{2mm}

\noindent\textbf{Q57. Feedback comes primarily from highly engaged users, while less engaged users rarely respond. Which approach best reduces participation bias?} \hfill [Tag: Feedback incorporation | Hard]
\begin{enumerate}[label=(\alph*),topsep=2pt,itemsep=1pt,partopsep=0pt,parsep=0pt]
    \item {Treat frequent respondents as representative}
    \item {Recruit targeted underrepresented users and compare patterns}
    \item {Use only the feedback with the strongest wording}
    \item {Discard all qualitative feedback}
\end{enumerate}
\vspace{2mm}

\noindent\textbf{Q58. A stakeholder requests a feature that benefits one department but increases delays for all others. Which evidence should guide the decision?} \hfill [Tag: Feedback incorporation | Hard]
\begin{enumerate}[label=(\alph*),topsep=2pt,itemsep=1pt,partopsep=0pt,parsep=0pt]
    \item {The feature's aggregate impact across affected workflows}
    \item {The number of words used in the request}
    \item {The feature's similarity to a competitor's product}
    \item {The requesting stakeholder's organizational seniority}
\end{enumerate}
\vspace{2mm}

\noindent\textbf{Q59. A continuous-improvement program repeatedly optimizes processing speed, but customer complaints increase. Which corrective principle is most appropriate?} \hfill [Tag: Continuous improvement of proposed solutions | Hard]
\begin{enumerate}[label=(\alph*),topsep=2pt,itemsep=1pt,partopsep=0pt,parsep=0pt]
    \item {Optimize the fastest workflow stage further}
    \item {Use a balanced set of quality and outcome measures}
    \item {Increase optimization frequency}
    \item {Remove complaint data because it is subjective}
\end{enumerate}
\vspace{2mm}

\noindent\textbf{Q60. A workflow's performance improves for several months, then declines as demand and user behavior change. Which improvement mechanism is most robust?} \hfill [Tag: Continuous improvement of proposed solutions | Hard]
\begin{enumerate}[label=(\alph*),topsep=2pt,itemsep=1pt,partopsep=0pt,parsep=0pt]
    \item {Repeat measurement, review assumptions, and adapt thresholds}
    \item {Remove variation from the performance data}
    \item {Declare the workflow permanently optimized}
    \item {Return automatically to the original workflow}
\end{enumerate}
\vspace{2mm}


\newpage
\section{Section B: Comprehensive Subjective Problem with Detailed Analytical Rubric}

\begin{subjectiveq}{5}{1}
\textbf{Case Background:} An enterprise software company is launching an automated AI workflow tool to summarize and draft complex legal contracts for corporate counsel. Initial customer interest is high, but pilot corporate customers express profound skepticism regarding generative hallucination and liability exposure.

\textbf{Required Design Deliverables:}
\begin{enumerate}[label=(\alph*)]
    \item Design a comprehensive Service Blueprint diagramming Frontstage, Backstage, Support Processes, and Failure Points across the contract review lifecycle.
    \item Architect a multi-fidelity prototyping and validation sequence (Smoke Test $\to$ Wizard of Oz $\to$ Functional Pilot).
    \item Formulate a Feedback Capture Grid and establish a 14-day rapid iterative refinement sprint plan.
\end{enumerate}
\end{subjectiveq}

\begin{solbox}
\textbf{1. Legal AI Service Blueprint Architecture:}\\
\begin{itemize}
    \item \textbf{Customer Actions:} Uploads 80-page non-disclosure and master services agreement; reviews annotated clause risk alerts; clicks 'Approve Clause'.
    \item \textbf{Frontstage Interactions:} Clean side-by-side contract viewer; color-coded risk flags (Red = high liability, Yellow = deviation from company playbook); confidence scores.
    \item \textbf{Line of Visibility} (What user cannot see):
    \item \textbf{Backstage Actions:} Document optical OCR parsing; semantic embeddings search against legal precedent database; multi-agent validation prompting.
    \item \textbf{Support Processes:} Secure isolated zero-data-retention server enclave; audit logging pipeline; attorney-in-the-loop verification routing.
    \item \textbf{Critical Failure Points:} Ingestion of non-standard scanned PDFs; missed indemnification sub-clauses; model latency exceeding 30 seconds.
\end{itemize}

\textbf{2. Validation \& Prototyping Progression:}\\
\begin{itemize}
    \item \textbf{Phase 1: Smoke Test / Fake Door:} Landing page with interactive demo video and 'Join Private Alpha for In-House Counsel' CTA button to validate market demand.
    \item \textbf{Phase 2: Wizard of Oz Prototype:} A seemingly automated web dashboard where uploads are processed behind the scenes by real human legal paralegals using standard playbooks. Validates user trust and interface layout without building full AI pipelines.
    \item \textbf{Phase 3: High-Fidelity Pilot with Human-in-the-Loop:} Real local LLM model deployed on-premise, with confidence thresholds: clauses below 95\% confidence require attorney click-through confirmation.
\end{itemize}

\textbf{3. Feedback Capture Grid \& 14-Day Sprint Loop:}\\
\begin{tabularx}{\linewidth}{|X|X|} \hline
\textbf{Likes (What resonated):} & \textbf{Criticisms (Friction points):} \\
$\bullet$ Instant identification of non-standard indemnification liability.\\
$\bullet$ Clean side-by-side diff view with suggested replacement text. &
$\bullet$ Obscure confidence percentages create legal confusion.\\
$\bullet$ Inability to export directly back to Microsoft Word tracked changes. \\ \hline
\textbf{Questions Asked:} & \textbf{New Ideas Generated:} \\
$\bullet$ 'Where is our sensitive contract data stored during parsing?'\\
$\bullet$ 'Can we fine-tune the engine on our proprietary company NDA?' &
$\bullet$ Instant 'Redline Export' directly to .docx with attorney notes.\\
$\bullet$ Custom toggle for Risk Appetite (Conservative vs Aggressive). \\ \hline
\end{tabularx}

\textit{Sprint Execution:} Days 1--3: Redesign confidence scores into clear qualitative categories. Days 4--8: Build .docx redline parser. Days 9--14: Re-test with 5 law firm pilot partners.
\end{solbox}

\begin{examinertip}
\textbf{Criteria for Full 10/10 Marks:}\\
1. \textbf{Rigorous Theoretical Framing (3 Marks):} Explicitly identify and define the formal model or framework applied (e.g., d.school stages, 5-Whys, MCDA, SCAMPER, Blueprinting, Minto SCQA). Avoid vague hand-waving.\\
2. \textbf{Contextual Application (4 Marks):} Directly connect every step to the specific facts, users, and constraints given in the case scenario.\\
3. \textbf{Actionable Deliverables \& Invariants (3 Marks):} Provide precise, formatted tables, matrices, or SMART criteria with clear quantitative bounds.
\end{examinertip}
